% texte.tex — f175v, vv. 1082–1127 (conversion initiale des travaux antérieurs)
\vers[n=1082, indent=3em, findent=1.7999999999999998em]{ΤΟ.}{Archer}{οὖτος τί λαλῖς;}{Holà ! qu'as-tu à bavarder ?}
\vers[suite]{ΕΥ.}{Euripide}{†οὖτος τί λαλῖς;†}{Holà ! qu'as-tu à bavarder ?}
\vers[n=1083, indent=3em, findent=1.7999999999999998em]{ΤΟ.}{Archer}{πρυτάνεις καλέσω;}{J'appellerai les prytanes.}
\vers[suite]{ΕΥ.}{Euripide}{†πρυτάνεις καλέσω;†}{J'appellerai les prytanes.}
\vers[n=1085, indent=3em, findent=1.7999999999999998em]{ΤΟ.}{Archer}{σί κακόν;}{Qu'est-ce qui arrive ?}
\vers[suite]{ΕΥ.}{Euripide}{†σί κακόν;†}{Qu'est-ce qui arrive ?}
\vers[n=1086, indent=3em, findent=1.7999999999999998em]{ΤΟ.}{Archer}{†πῶτε† τὸ πωνή;}{D'où part cette voix ?}
\vers[suite]{ΕΥ.}{Euripide}{†πῶτε τὸ πωνή;†}{D'où part cette voix ?}
\vers[n=1087, indent=3em, findent=1.7999999999999998em]{ΤΟ.}{Archer}{σὺ λαλῖς;}{Est-ce toi qui bavardes?}
\vers[suite]{ΕΥ.}{Euripide}{†σὺ λαλῖς;†}{Est-ce toi qui bavardes?}
\vers[indent=3em, findent=1.7999999999999998em]{ΤΟ.}{Archer}{κλαύσαι.}{Gare à toi.}
\vers[suite]{ΕΥ.}{Euripide}{†κλαύσαι.†}{Gare à toi.}
\vers[n=1088, indent=3em, findent=1.7999999999999998em]{ΤΟ.}{Archer}{κἀκκάσκι μοι;}{Tu te gausses de moi.}
\vers[suite]{ΕΥ.}{Euripide}{†κἀκκάσκι μοι;†}{Tu te gausses de moi.}
\vers[n=1090, indent=3em, findent=1.7999999999999998em]{ΧΟ.}{Chœur}{μὰ Δί᾽ ἀλλὰ γυνὴ πλησίον αὕτη.}{Non, par Zeus : c'est cette femme, là tout près.}
\vers[n=1091, indent=3em, findent=1.7999999999999998em]{ΕΥ.}{Euripide}{πλήξασ᾿ αὕτη.}{Là tout près.}
\vers[n=1092, indent=3em, findent=1.7999999999999998em]{ΤΟ.}{Archer}{ποῦ ᾿στ᾽ ἠ μιαρά;}{Où est la canaille ?}
\vers[suite]{ΧΟ.}{Chœur}{καὶ δὴ φεύγει.}{La voilà qui se sauve.}
\vers[n=1093, indent=3em, findent=1.7999999999999998em]{ΤΟ.}{Archer}{ποῖ ποῖ πεύγεις;}{Où, où te sauves-tu ?}
\vers[suite]{ΕΥ.}{Euripide}{ποῖ ποῖ πέτει.}{Où, où voles-tu ?}
\vers[n=1094, indent=3em, findent=1.7999999999999998em]{ΤΟ.}{Archer}{οὐ καιρήσεις.}{Il t'en cuira.}
\vers[suite]{ΕΥ.}{Euripide}{οὐχ αἱρήσεις.}{Tu ne m'auras.}
\vers[n=1095, indent=3em, findent=1.7999999999999998em]{ΤΟ.}{Archer}{ἔτι γὰρ γρύζεις;}{Tu maronnes encore ?}
\vers[suite]{ΕΥ.}{Euripide}{†ἔτι γὰρ γρύζεις;†}{Tu maronnes encore ?}
\vers[n=1096, indent=3em, findent=1.7999999999999998em]{ΤΟ.}{Archer}{λαβὲ τὴ μιαρά.}{Arrête, la canaille.}
\vers[suite]{ΕΥ.}{Euripide}{†λαβὲ τὴ μιαρά.†}{Arrête, la canaille.}
\vers[n=1097, indent=3em, findent=1.7999999999999998em]{ΤΟ.}{Archer}{λάλο καὶ κατάρατο γύναικο.}{La bavarde, la maudite femme.}
\vers[n=1098]{ΕΥ.}{Euripide}{ὦ θεοὶ τίν᾽ εἰς γῆν βαρβάρων ἀφίγμεθα}{\didasc{En Persée :} « En quel pays barbare, ô dieux, ces talonnières}
\vers[n=1099]{}{}{ταχεῖ πεδίλῳ; διὰ μέσου γὰρ αἰθέρος}{ont-elles amené Persée ? Au milieu de l'éther}
\vers[n=1100]{}{}{τέμνων κέλευθον πόδα τίθημ᾽ ὑπόπτερον}{me frayant un chemin, prestement je fends l'air d'un pied ailé,}
\vers[n=1101]{}{}{Περσεὺς πρὸς Ἄργος ναυστολῶν τὸ Γοργόνος}{portant à la cité argienne la tête de}
\vers[n=1102]{}{}{κάρα κομίζων.}{Gorgone. »}
\vers[suite]{ΤΟ.}{Archer}{τί λέγι; τὴ Γόργος πέρι}{Que dis-tu? Tu portes la tête de Gorgon}
\vers[n=1103]{}{}{τὸ γραμματέο σὺ τὴ κεπαλή;}{le scribe ?}
\vers[suite]{ΕΥ.}{Euripide}{τὴν Γοργόνος}{De Gorgone, ai-je dit.}
\vers[n=1104]{}{}{ἔγωγε φημί.}{ἔγωγε φημί.}
\vers[suite]{ΤΟ.}{Archer}{Γόργο τοι κἀγὼ λέγι.}{Gorgon, parfaitement, je dis de même.}
\vers[n=1105]{ΕΥ.}{Euripide}{ἔα· τίν᾽ ὄχθον τόνδ᾽ ὁρῶ καὶ παρθένον}{« Ciel! Quelle est cette vierge aux déesses pareille amarrée à ce roc,}
\vers[n=1106]{}{}{θεαῖς ὁμοίαν ναῦν ὅπως ὡρμισμένην;}{comme un navire au port ? »}
\vers[n=1107]{ΜΝ.}{Mnésiloque}{ὦ ξένε κατοίκτιρόν με τὴν παναθλίαν,}{« Étranger, prends pitié de moi, la désolée.}
\vers[n=1108]{}{}{λῦσόν με δεσμῶν.}{Détache mes liens. »}
\vers[suite]{ΤΟ.}{Archer}{οὐκὶ μὶ λαλῆσι σύ;}{Veux-tu ne pas bavarder ?}
\vers[n=1109]{}{}{κατάρατο τολμᾷς ἀποτανουμένη λαλῖς;}{Maudit, oses-tu bien ? Tu vas mourir et tu bavardes.}
\vers[n=1110]{ΕΥ.}{Euripide}{ὦ παρθέν᾽ οἰκτίρω σὲ κρεμαμένην ὁρῶν.}{« Ô vierge, une pitié me prend quand je te vois suspendue à ce roc. »}
\vers[n=1111]{ΤΟ.}{Archer}{οὐ παρτέν᾽ ἐστίν, ἀλλ᾽ ἀμαρτωλὴ γέρων}{Ce n'est pas une vierge, c'est un vieux coquin,}
\vers[n=1112]{}{}{καὶ κλέπτο καὶ πανοῦργο.}{et un voleur, et un fripon.}
\vers[suite]{ΕΥ.}{Euripide}{ληρεῖς ὦ Σκύθα.}{Tu plaisantes, ô Scythe.}
\vers[n=1113]{}{}{αὕτη γάρ ἐστιν Ἀνδρομέδα παῖς Κηφέως.}{C'est Andromède, fille de Cépheus.}
\vers[n=1114]{ΤΟ.}{Archer}{σκέψαι τὸ κύστο· μή τι μικκὸν παίνεται;}{Regarde son cas : te paraît-il petit ?}
\vers[n=1115]{ΕΥ.}{Euripide}{φέρε δεῦρό μοι τὴν χεῖρ᾽, ἵν᾽ ἅψωμαι κόρης·}{Donne-moi la main, que j'atteigne à la jeune fille.}
\vers[n=1116]{}{}{φέρε Σκύθ᾽· ἀνθρώποισι γὰρ νοσήματα}{Voyons, Scythe « chaque homme a ses faiblesses » :}
\vers[n=1117]{}{}{ἅπασίν ἐστιν· ἐμὲ δὲ καὐτὸν τῆς κόρης}{la mienne est d'être féru}
\vers[n=1118]{}{}{ταύτης ἔρως εἴληφεν.}{d'amour pour cette jeune fille.}
\vers[suite]{ΤΟ.}{Archer}{οὐ ζηλῶσί σε·}{Je ne suis pas jaloux.}
\vers[n=1119]{}{}{ἀτὰρ εἰ τὸ πρωκτὸ δεῦρο περιεστραμμένον,}{S'il tournait les fesses par ici,}
\vers[n=1120]{}{}{οὐκ ἐπτόνησά σ᾽ αὐτὸ πυγίζεις ἄγων.}{je ne te refuserais pas d'en passer ta fantaisie.}
\vers[n=1121]{ΕΥ.}{Euripide}{τί δ᾽ οὐκ ἐᾷς λύσαντά μ᾽ αὐτὴν ὦ Σκύθα}{Que ne me permets-tu de la délier, ô Scythe}
\vers[n=1122]{}{}{πεσεῖν ἐς εὐνὴν καὶ γαμήλιον λέχος;}{« et de me laisser choir dans le lit nuptial »}
\vers[n=1123]{ΤΟ.}{Archer}{εἰ σπόδρ᾽ ἐπιτυμεῖς τὴ γέροντο πύγισο,}{Si tu as tant envie de jouir de ce vieillard,}
\vers[n=1124]{}{}{τὴ σανίδο τρήσας ἐξόπιστο πρώκτισον.}{perce-moi l'ais par derrière et enfourne.}
\vers[n=1125]{ΕΥ.}{Euripide}{μὰ Δί᾽ ἀλλὰ λύσω δεσμά.}{Non, par Zeus, je briserai ses liens.}
\vers[suite]{ΤΟ.}{Archer}{μαστιγῶ σ᾽ ἄρα.}{Alors je te fouette.}
\vers[n=1126]{ΕΥ.}{Euripide}{καὶ μὴν ποιήσω τοῦτο.}{N'importe, je vais le faire.}
\vers[suite]{ΤΟ.}{Archer}{τὸ κεπαλή σ᾽ ἄρα}{Alors je te coupe la tête}
